% Source: https://artofproblemsolving.com/community/c3h1384140_a_really_nice_nt_problem
% Stable identity: aops-topic-1384140-q01
\textbf{(a)} Take $r=2$, $k=4$. Then $rk=8$ has $y=4$ divisors, but the only pairs with $\gcd=2$ and $\operatorname{lcm}=4$ are $(2,4)$ and $(4,2)$. The argument fails because $mn=rk$ is only a necessary condition: most factorisations $rk=mn$ (here $1\cdot8$, $8\cdot1$) do not have the right gcd and lcm.

\textbf{(b)} If $r\nmid k$ there are no pairs, since $\gcd(m,n)$ divides $\operatorname{lcm}(m,n)$. Suppose $r\mid k$ and write $m=ra$, $n=rb$. Then $\gcd(m,n)=r$ means $\gcd(a,b)=1$, and then $\operatorname{lcm}(m,n)=rab$, so we need
\[
ab=\dfrac {k}{r},\qquad \gcd(a,b)=1 .
\]
If $\dfrac {k}{r}=p_1^{e_1}\cdots p_t^{e_t}$, coprimality forces each whole prime power $p_i^{e_i}$ to go entirely into $a$ or entirely into $b$, and every such choice works. So the number of ordered pairs is
\[
\boxed{2^{t}},\qquad t=\text{the number of distinct primes dividing }\dfrac {k}{r}
\]
(and $0$ if $r\nmid k$). In the example, $\dfrac {k}{r}=2$, $t=1$, giving $2$ pairs.
